\documentclass[a4paper,11pt,fleqn]{article}%fleqn pour aligner les equations à gauche
\usepackage{../../Styles/Eval}


\newcommand{\chnum}{Chapitre 1 et 2}
\newcommand{\chtitle}{Transformation acide-base et pH}
\newcommand{\dtype}{DS 1}

\begin{document}
\maketitle

\section{Beurre acide}
\noindent \colorbox{red!10}{\begin{minipage}{.98\linewidth}
	\small
	\textbf{Acide butyrique}\\
	Les travaux de Crevreul lui ont permis d’isoler un acide gras à quatre atomes de carbone, qu’il a nommé « acide butyrique », dérivé du latin \textit{butyrum} signifiant « \textit{beurre} », à cause de son odeur caractéristique de beurre rance. Pour la suite, l’élément but(y)- a été repris pour former les noms de molécules à quatre atomes de carbone dont, avec le suffixe -ane des alcane, celui du butane.\\
\end{minipage}}

L’acide butyrique est un acide carboxylique possédant une chaîne carbonée, dérivée du butane. Il est présent dans le couple acide butyrique/ion butyrate (\ce{C4H8O2}/\ce{C4H7O2-}).
\begin{enumerate}
	\item Déterminer le schéma de Lewis de l’acide butyrique. (1pt)
	\item Justifier le caractère "acide" de la molécule. (1pt)
	\item Écrire la demi-équation associée à l'acide butyrique. (1pt)
	\item Écrire l'équation de la réaction mise en jeu lorsque l'acide butyrique est mis dans l'eau. (1pt)
	\item On réalise une solution d'acide butyrique à partir d'une masse de \SI{2,2}{\milli\gram} dissous dans \SI{500}{\milli\liter} d'eau. En partant de l'hypothèse que tout l'acide réagit avec l'eau, déterminer le pH de la solution obtenue. (1pt)
\end{enumerate}
Le butanoate de méthyle \ce{CH3(CH2)2COOCH3} est un composé chimique utilisé comme arôme pour son odeur de pomme. Il est obtenu à partir de l’acide butyrique.
\begin{enumerate}[resume]
	\item Recopier puis compléter le schéma de Lewis du butanoate de méthyle suivant, puis identifier les éventuels groupes fonctionnels. (1pt)
	\item Préciser si l’ester formé a un caractère acide. (1pt)
	\item Attribuer chacun des spectres suivants à la molécule (acide butyrique ou butanoate de méthyle) correspondante en justifiant. (1pt)
\end{enumerate}

\begin{figure}[h!]
	\centering \chemfig[atom sep = 2em]{CH_3-CH_2-CH_2-C(=[1]O)-[7]O-CH_3}
	\caption{Butanoate de méthyle}

	\begin{minipage}{.48\linewidth}
		\includegraphics[width=.9\linewidth]{Ressources/Butanoate de méthyle1.png}
	\end{minipage}\hspace{.02\linewidth}
	\begin{minipage}{.48\linewidth}
		\includegraphics[width=.9\linewidth]{Ressources/Butanoic1.png}
	\end{minipage}
	\caption{Spectres IR à attribuer}
\end{figure}

\newpage

\section{Coloration capillaire}
Les coloration capillaires classiques sont des solutions basiques à base d'ammoniaque. Elles sont classées en deux catégories en fonction de l'utilisation : les colorations commerciales, avec une teneur de \num{27} à \SI{30}{\percent} d'ammoniaque et les mélanges professionnels avec une teneur maximale de \SI{6}{\percent}.\\

On propose de déterminer la concentration en ammoniaque (\ce{NH3}) d'une solution servant de base pour une préparation commerciale de coloration capillaire notée $S_0$ de concentration notée $c_0$. Pour cela, on réalise un titrage sur \SI{10}{mL} de solution diluée,  grâce à une solution titrante d'acide chlorhydrique \ce{HCl} de concentration $c_A$ = \SI{1E-3}{\mole\per\liter}.

Le suivi de la réaction se fait grâce à une sonde de conductimétrie et les points mesurés sont reportés dans le graphique du document fourni.

	\subsection{Étude préliminaire}

\begin{enumerate}
	\item La solution à titrer semble trop concentrée pour un titrage direct, une  d'un facteur 10 est donc réalisé pour créer une solution titrée $S_1$ de concentration notée $c_1$. Lister le matériel nécessaire pour réaliser \SI{50}{mL} de solution $S_1$. (1pt)
	\item Quelles sont les limites d'un suivi par conductimétrie ? (1pt)
	\item Écrire la réaction de se produisant entre l'ammoniac \ce{NH3} et l'eau. Puis celle se produisant entre l'acide chlorhydrique \ce{HCl} et l'eau. (2pts)
	\item En déduire la réaction support du titrage. (1pt)
	\item Réaliser un schéma du titrage. (1pt)
	\item Compléter le tableau suivant en faisant figurer l'évolution de la quantité de matière de chaque espèce ionique dans le milieu au cours du titrage.(1pt)
\end{enumerate}
\begin{table}[h!]
\centering
	\begin{tabular}{|>{\centering}m{0.2\textwidth}|>{\centering}m{0.15\textwidth}|>{\centering}m{0.15\textwidth}|>{\centering}m{0.15\textwidth}|>{\centering\arraybackslash}m{0.15\textwidth}|} \hline
	\textbf{Espèce ionique} & &&&\\ \hline
	\textbf{Avant équivalence} & &&&\\ \hline
	\textbf{A l'équivalence} & &&&\\ \hline
	\textbf{Après l'équivalence} & &&&\\ \hline
	\end{tabular}
\end{table}

	\subsection{Exploitation des résultats}

\begin{enumerate}[resume]
	\item En vous aidant du tableau précédent, et des valeurs de conductivité molaires ioniques fournies, justifier l'allure de la courbe de titrage. (2pts)
	\item Déterminer, en explicitant la méthode employée, la valeur du volume équivalent. (1pt)
	\item Calculer la valeur de la concentration $c_1$. (2pts)
	\item Déterminer la valeur de la concentration de la solution de coloration capillaire. (0.5pt)
	\item Vérifier que les valeurs obtenues sont en accord avec la teneur en ammoniaque fournie par le vendeur.(0.5pt)
\end{enumerate}

\begin{figure}[h!]
\centering
	\includegraphics[width = .7\textwidth]{Ressources/TSPE_DS2_Fig2.png}
    \caption{Titrage de la solution $S_1$ par l'acide chlorhydrique.}
\end{figure}

\begin{minipage}{.97\textwidth}
\textbf{Données : }
\begin{itemize}
	\item \textbf{Couples acide/base :} \ce{NH4+}/\ce{NH3} ; \ce{H3O+}/\ce{H2O} ; \ce{H2O}/\ce{HO-} 
	\item \textbf{Densité de la base pour coloration capillaire : }$d$ = \num{0.9}
	\item \textbf{Masse molaire de l'ammoniac : } $M$ = \SI{17}{\gram\per\mole}
	\item \textbf{Conductivités molaires ioniques : }\\
	$\lambda$(\ce{NH4+}) = \SI{7.35}{\milli\siemens\meter\squared\per\mole}\\
	$\lambda$(\ce{Cl-}) = \SI{7.63}{\milli\siemens\meter\squared\per\mole}\\
	$\lambda$(\ce{HO-}) = \SI{19.8}{\milli\siemens\meter\squared\per\mole}\\
	$\lambda$(\ce{H3O+}) = \SI{35.0}{\milli\siemens\meter\squared\per\mole}\\
\end{itemize}
\end{minipage}

\textbf{Tableau des références IR}\\

\begin{tabular}{|>{\centering}m{.15\linewidth}|>{\centering}m{.45\linewidth}|>{\centering \arraybackslash} m{.3\linewidth}|}
    \hline
    \textbf{Liaison} &\textbf{Condition} & \textbf{Nombre d'onde} $\overline{\nu}$ (\unit{\per\centi\meter})\\
    \hline
    \chemfig{O-H} & Sans pont hydrogène & \num{3580} - \num{3650} (bande fine)\\ \cline{2-3}
    & Avec pont hydrogène & \num{3200} - \num{3400} (bande large)\\ \cline{2-3}
    & Si cette liaison  fait partie du groupe caractéristique carbonyle & \num{2500} - \num{3200} (bande large)\\ \cline{2-3}
    %\hline
    %\chemfig{N-H} & - & \num{3400} - \num{3500}\\
    %\hline
    %\chemfig{C-H} & - & \num{2800} - \num{3100}\\ \cline{2-3}
    %& Si l'atome de carbone est lié par une liaison double avec un atome d'oxygène & \num{2750} - \num{2900}\\ \cline{2-3}
    \hline
    \chemfig{C=O} & Si cette liaison fait partie du groupe carbonyle & \num{1650} - \num{1730}\\ \cline{2-3}
    & Si cette liaison fait partie du groupe carboxyle & \num{1680} - \num{1710}\\ \cline{2-3}
    & Si cette liaison fait partie du groupe caractéristique ester & \num{1700} - \num{1740}\\ \cline{2-3}
    & Si cette liaison fait partie du groupe caractéristique amide & \num{1650}\\
    \hline
    \chemfig{C=C} & Si cette liaison ne fait pas partie d'un sycle aromatique & \num{1640} - \num{1680}\\ \cline{2-3}
    & Si cette liaison fait partie d'un sycle aromatique & \num{1500} - \num{1580}\\
    \hline 
\end{tabular}
\end{document}